\documentclass[11pt]{article}
\usepackage[utf8]{inputenc}
\usepackage[spanish]{babel}
\usepackage{amsmath,amssymb,amsthm}
\usepackage{graphicx}
\usepackage[margin=2.5cm]{geometry}
\usepackage{enumitem}

\title{Semana 3 -- Ejercicios}
\author{Tutorías -- Introducción al Cálculo}
\date{17/08--21/08/2026}

\begin{document}
\maketitle

\section{Ejercicios de la semana}

\begin{enumerate}[label=\arabic*.]
  \item Usando los axiomas de orden y el Teorema 2.3, demuestre que si $x<0$ y $y<0$, entonces
  $x\cdot y>0$. (Indicación: use (III) y (VII) del Teorema 2.3.)

  \item Demuestre que si $x\neq 0$, entonces $x^2>0$. (Indicación: combine (X) del Teorema 2.3 con
  (III) del Teorema 2.4.)

  \item Sean $a,b\in\mathbb{R}$ con $a<b$. Demuestre que no existe ningún $x\in\mathbb{R}$ tal que
  $x<a$ y $x>b$ simultáneamente.

  \item Calcule, usando la definición, el valor de $|-4|+|4|-|{-3}\cdot 2|$. (Respuesta: $2$.)

  \item Determine todos los $x\in\mathbb{R}$ tales que $|2x-1|=5$. (Indicación: use (IV) del
  Teorema 2.4.) (Respuesta: $x=3 \lor x=-2$.)

  \item Determine todos los $x\in\mathbb{R}$ tales que $|x+2| \le 3$. (Respuesta: $[-5,1]$.)

  \item Demuestre que para todo $x\in\mathbb{R}$, $|x|^2=x^2$.

  \item Demuestre, usando la desigualdad triangular, que para todo $x,y,z\in\mathbb{R}$ se cumple
  $|x-z| \le |x-y|+|y-z|$.

  \item Sean $x,y\in\mathbb{R}$ positivos. Demuestre que $|x-y| < x+y$.

  \item \textbf{(Nivel interrogación)} Demuestre rigurosamente, usando \emph{solo} los axiomas de
  orden (multiplicando la desigualdad $a<b$ primero por $a$ y luego por $b$, y usando
  transitividad), que para todo $a,b\in\mathbb{R}$ con $0<a<b$ se tiene $a^2<b^2$. Indique
  explícitamente qué axioma utiliza en cada paso.

  \item \textbf{(Nivel interrogación)} Resuelva la ecuación $|x^2-4| = 3x$. (Indicación: separe en
  casos según el signo de $x^2-4$, y recuerde que el lado derecho debe ser mayor o igual a $0$
  para que la igualdad sea posible.)

  \item \textbf{(Nivel interrogación)} Demuestre rigurosamente, usando solo los axiomas de orden,
  que si $a<b$ y $c<0$, entonces $ac>bc$.
\end{enumerate}

\end{document}
